\subsection{An actual reduction at extreme corners}

\begin{lemma}[Proved: an extreme one-row corner reduces to its minimal group]
\label{CP2:A-extreme-minimal-group-reduction}
Let $x$ be a source object in a finite category $C$: no other
object has an arrow to $x$. Write $M=C(x,x)$, choose an
idempotent $e$ in its minimal ideal, and put $G=eMe$.
If $Me=G$, there is a finite category $C'$ of the same
homotopy type obtained by replacing the corner $M$ by $G$
and every outgoing stalk $C(x,y)$ by $C(x,y)e$. In particular,
\[
 |\operatorname{Mor}C'|\leq
 |\operatorname{Mor}C|-\bigl(|M|-|G|\bigr).
\]
The hypothesis $Me=G$ is equivalent to the minimal ideal
having one row. There is a dual statement for a sink object
with $eM=G$, equivalently one minimal column.
\end{lemma}
\begin{proof}
Adjoin the retract object $e$ to $C$, inside its idempotent
completion, and call the resulting category $C^e$. The old
object $x$ retracts onto the added object, so the inclusion
$BC\to BC^e$ is a homotopy equivalence by the retract-object
lemma. Let $C'$ be the full category obtained from $C^e$
by deleting $x$.

Apply Quillen's Theorem~A to the inclusion $i:C'\to C^e$,
using the comma categories $(i\downarrow z)$. If $z$ is
retained, its identity is terminal in the comma category.
For $z=x$, no old retained object has an arrow into $x$.
Thus the comma category is precisely the right action
category of $G$ on
\[
 C^e(e,x)=Me=G.
\]
Right multiplication is free and transitive. Consequently
this comma category has exactly one morphism between each
ordered pair of objects and is contractible. The theorem
gives $BC'\simeq BC^e\simeq BC$.

The retained old hom-sets are unchanged, the new corner is
$G$, and the new outgoing hom-set to $y$ is $C(x,y)e$.
These are subsets of the corresponding old hom-sets, and
there are no new incoming arrows. This proves the arrow
inequality. The finite minimal-ideal description identifies
$Me$ with a disjoint union of one free transitive $G$-set
for each minimal row, proving the stated equivalence of
hypotheses. Opposition proves the sink assertion.
\end{proof}

\begin{corollary}[Proved: the extreme-corner cut at the next arrow budget]
\label{CP2:A-extreme-minimal-group-eighteen}
An eighteen-arrow category of type $\CP^2$ cannot have
a source corner which is a nongroup monoid with one minimal
row, or a sink corner which is a nongroup monoid with one
minimal column.
\end{corollary}
\begin{proof}
For a nongroup monoid, its minimal group has strictly fewer
elements: equality would imply $M=G$ as a set and hence as
a monoid. The preceding actual homotopy equivalence would
give a category with at most seventeen arrows, contrary
to Theorem~\ref{CP2:A-global-lower-eighteen}.
\end{proof}

The literal table predicate uses only minimal-ideal data.
For an idempotent $e$ in a corner's minimal ideal it computes
the sets $Me$, $eM$, and $eMe$ directly. At a source, the
test is $|Me|=|eMe|<|M|$; at a sink it is
$|eM|=|eMe|<|M|$. Since $eMe$ is contained in both other
sets, equality of cardinalities is equality of sets.
This test involves neither a homology approximation nor
an assumption that the chosen corner action is trivial.
